\documentclass[11pt]{article}

\usepackage{series} % user-provided
\usepackage{amsmath,amssymb,amsthm,amsfonts}
\usepackage{mathtools}
\usepackage{enumitem}

\hypersetup{colorlinks=true, linkcolor=blue, citecolor=blue, urlcolor=blue}

% ============================
% Theorem environments
% ============================
\theoremstyle{plain}

\theoremstyle{definition}

% ============================
% Notation
% ============================
\newcommand{\Spec}{\mathrm{Spec}}

\title{Black Hole Information Loss and Quantum Measurement Collapse\\
as the Same Admissibility Transition in Modal Triplet Theory}
\author{Peter Nero}
\date{January, 2026}

\begin{document}
\maketitle

\begin{abstract}
We propose and formalize a single mechanism in Modal Triplet Theory (MTT) that unifies two
apparently unrelated 4D problems: (i) black hole information loss in semiclassical gravity and
(ii) quantum measurement collapse in nonrelativistic quantum mechanics. In MTT, 4D physics is
the shadow of deterministic dynamics on a higher configuration space under a noninvertible
coherent projection. We define admissibility barriers as loci where the coherent projection fails
to be stably invertible (spectral gap closure / projector discontinuity / basin rearrangement).
We prove a general ``projection noninvertibility'' theorem: whenever trajectories cross an
admissibility barrier, the induced 4D evolution becomes noninvertible and admits no global
information recovery map, despite invertibility of the upstairs evolution. We then show how
measurement selection and horizon formation/evaporation are two instantiations of barrier
crossing, and we map the resulting 4D shadows to (a) collapse/Born probabilities and (b)
thermal Hawking radiation/mixed exterior states. Finally, we relate the framework to the Page
curve and island formula program: islands are interpreted as regime-dependent partial inversions
(or re-encodings) of the projection on a restricted algebra, rather than restoration of global
4D invertibility.
\end{abstract}

\tableofcontents

% ======================================================================
\section{Scope and Claim Discipline}
% ======================================================================

This paper has two goals:
\begin{enumerate}[label=(G\arabic*)]
\item Provide a theorem-level mechanism in MTT that yields loss of invertibility in the 4D
shadow dynamics whenever trajectories cross an admissibility barrier.
\item Show that both (i) quantum measurement collapse and (ii) black hole information loss are
instances of that same mechanism under different physical identifications of the barrier.
\end{enumerate}

We separate three layers:
\begin{enumerate}[label=(L\arabic*)]
\item \textbf{Mathematical layer (proved):} deterministic evolution on a configuration space,
a noninvertible projection, and barrier-crossing implies noninvertibility of the projected
dynamics.
\item \textbf{Interpretive layer (identified):} measurement events correspond to constraint/context
changes that move the system across an admissibility barrier; horizon formation/evaporation
corresponds to a barrier between exterior-admissible and interior/nonadmissible sectors.
\item \textbf{Phenomenological layer (compared):} the resulting shadows align with (a) collapse
and Born probabilities and (b) thermal radiation / mixed exterior states, and clarify the role of
islands/replica-wormhole ``patches'' as partial re-encodings rather than global invertibility
restoration.
\end{enumerate}

We do \emph{not} assume fundamental 4D nonunitarity. The only noninvertibility asserted is at
the level of the projected 4D shadow map.

% ======================================================================
\section{Upstairs Setup: Dynamics, Projection, and Admissibility Barriers}
% ======================================================================

\subsection{Configuration space and deterministic evolution}

Let $(\mathcal X,\mathcal B)$ be a standard Borel space representing the full MTT configuration
space on a bounded-geometry slab (fields on $M_{10}$ with suitable Sobolev topology, restricted
to a finite-energy domain). Let
\[
\Phi_t : \mathcal X \to \mathcal X
\]
be a measurable, invertible flow (or discrete-time map $\Phi:\mathcal X\to\mathcal X$) representing
the deterministic upstairs evolution. Invertible means $\Phi_t$ is bijective with measurable
inverse.

\begin{remark}
In MTT the upstairs dynamics is a well-posed flow on the full modal system. The present section
requires only invertibility and measurability; no specific PDE form is needed.
\end{remark}

\subsection{Coherent projection and observable map}

Let $\Pi:\mathcal X\to \mathcal X_{\mathrm{coh}}$ be a measurable projection to the coherent sector,
and let $I:\mathcal X_{\mathrm{coh}}\to \mathcal Y$ be an ``observable pushforward'' to an effective 4D
state space $\mathcal Y$ (e.g.\ a Hilbert-space state space or state on a net of algebras). Define
the 4D observable map
\[
P := I\circ \Pi : \mathcal X \to \mathcal Y.
\]

\begin{definition}[Projected (shadow) dynamics]
Given an initial condition $x_0\in\mathcal X$, define the shadow trajectory
\[
y(t) := P(\Phi_t(x_0)) \in \mathcal Y.
\]
\end{definition}

\subsection{Admissibility and barrier sets}

MTT asserts that coherent-sector evolution is only physically asserted on an admissible subset
$\mathcal A\subset \mathcal X$ (bounded geometry, spectral gap persists, projector boundedness,
FCC stability margins positive, etc.). We encode this abstractly:

\begin{definition}[Admissible set]
An admissible set $\mathcal A\subset \mathcal X$ is a measurable subset such that: for initial data
in $\mathcal A$, the coherent projection $\Pi$ is well-defined, bounded/regular in the relevant
topologies, and stable under perturbations, and the induced coherent dynamics remains in a
well-controlled regime for slab-local times.
\end{definition}

The key new concept is the barrier:

\begin{definition}[Admissibility barrier]
A measurable set $\mathcal B\subset \mathcal X$ is an admissibility barrier if:
\begin{enumerate}[label=(B\arabic*)]
\item $\mathcal X\setminus \mathcal B$ splits into at least two measurable regions
$\mathcal U_+$ and $\mathcal U_-$ such that $\mathcal U_+\cap\mathcal U_-=\emptyset$,
$\mathcal U_+\cup\mathcal U_-=\mathcal X\setminus\mathcal B$.
\item $\Pi$ restricted to each $\mathcal U_\pm$ is regular (continuous in the chosen topology,
or at minimum locally constant on fibers) but there is no globally measurable right-inverse
for $\Pi$ across $\mathcal U_+\cup\mathcal U_-$.
\item Basin/sector structure changes across $\mathcal B$ in the sense that for some
$y\in\mathcal Y$, the fiber $P^{-1}(y)$ intersects both $\mathcal U_+$ and $\mathcal U_-$.
\end{enumerate}
\end{definition}

\begin{remark}[Physical interpretation of barriers]
In MTT, barriers arise when spectral gaps close, coherent projectors lose regularity, or stability
margins vanish. The definition above is the measure-theoretic abstraction of those events.
\end{remark}

% ======================================================================
\section{The Bridge Theorem: Barrier Crossing Implies Shadow Noninvertibility}
% ======================================================================

We now state the core theorem: invertibility upstairs does not descend through a noninvertible
projection; barrier crossing produces irrecoverable loss of 4D reconstructability.

\subsection{Noninvertibility of the projected evolution}

\begin{definition}[Shadow evolution operator]
Define the shadow map at time $t$ by
\[
T_t := P\circ \Phi_t : \mathcal X \to \mathcal Y.
\]
\end{definition}

\begin{definition}[Shadow invertibility]
We say the shadow at time $t$ is (globally) invertible if there exists a measurable map
$S_t:\mathcal Y\to \mathcal X$ such that
\[
T_t\circ S_t = \Id_{\mathcal Y}.
\]
(That is, $S_t$ is a measurable right-inverse of $T_t$.)
\end{definition}

\begin{theorem}[Barrier crossing implies noninvertible shadow dynamics]
\label{thm:barrier-noninvertibility}
Assume:
\begin{enumerate}[label=(H\arabic*)]
\item $\Phi_t$ is invertible and measurable on $\mathcal X$.
\item $P:\mathcal X\to\mathcal Y$ is measurable.
\item $\mathcal B\subset\mathcal X$ is an admissibility barrier with regions $\mathcal U_\pm$ as in
Definition 2.3.
\end{enumerate}
Then for any time $t$ such that $\Phi_t(\mathcal U_+)\cap \Phi_t(\mathcal U_-)\neq\emptyset$ and
$P^{-1}(y)$ intersects both regions for some $y\in\mathcal Y$, the shadow map $T_t$ admits no
measurable right-inverse; in particular, shadow dynamics is not globally invertible.
\end{theorem}

\begin{proof}
By Definition 2.3(B3), there exists $y\in\mathcal Y$ and points $x_+\in\mathcal U_+$,
$x_-\in\mathcal U_-$ such that $P(x_+)=P(x_-)=y$. Since $\Phi_t$ is invertible, the images
$\Phi_t(x_+),\Phi_t(x_-)$ are distinct.

Compute:
\[
T_t(x_+) = P(\Phi_t(x_+)),\qquad T_t(x_-) = P(\Phi_t(x_-)).
\]
If $T_t$ had a measurable right-inverse $S_t$, then $x_+=S_t(T_t(x_+))$ and
$x_-=S_t(T_t(x_-))$ would hold for all inputs, in particular on points with the same image.
But since $P$ is noninjective on fibers crossing the barrier, there exist points with identical
shadow $y$ but distinct upstairs origins, and no single-valued right-inverse can select both in a
measurable way on a set of positive measure without violating $T_t\circ S_t=\Id_{\mathcal Y}$.
Hence no measurable right-inverse exists.
\end{proof}

\begin{corollary}[Upstairs unitarity does not imply 4D unitarity]
\label{cor:unitarity}
Even when $\Phi_t$ is invertible/unitary at the upstairs level, the induced 4D shadow dynamics
cannot be globally invertible once barrier-crossing noninjectivity occurs. Any 4D ``unitarity
restoration'' must therefore rely on enlarging $\mathcal Y$ or introducing nonlocal encodings.
\end{corollary}

\begin{remark}[Why this is the common core of both paradoxes]
Both measurement collapse and black hole information loss are precisely questions of whether
the 4D map $T_t$ can be made invertible/unitary when the underlying theory is invertible.
Theorem~\ref{thm:barrier-noninvertibility} answers: not if barrier crossing produces fiber
identifications.
\end{remark}

% ======================================================================
\section{Two Shadows of One Barrier: Measurement and Black Holes}
% ======================================================================

\subsection{Measurement collapse as barrier crossing}

In the measurement setting, a ``context'' $C$ (apparatus coupling, POVM/instrument, etc.)
induces an admissible set $\mathcal A_C\subset\mathcal X$ and an associated basin decomposition
into outcome-attractors. A measurement event corresponds to a change of context
$C\to C'$ that changes admissibility constraints and can move trajectories across an
admissibility barrier $\mathcal B$ separating basin atlases.

\begin{proposition}[Collapse as projected noninvertibility]
In the measurement setting, outcome selection corresponds to barrier crossing in $\mathcal X$,
and the impossibility of reconstructing the pre-measurement coherent configuration from the
post-measurement outcome is an instance of Theorem~\ref{thm:barrier-noninvertibility}.
\end{proposition}

\begin{remark}
In MTT, probabilities arise from basin measures/weights on the admissible set. This paper does
not re-derive the Born rule; it uses only the structural fact that basin selection is projection-
induced and noninvertible in the 4D shadow.
\end{remark}

\subsection{Black hole information loss as barrier crossing}

In the black hole setting, consider a collapse+evaporation history on a slab that includes horizon
formation. The relevant claim is not that upstairs information is destroyed, but that the 4D
shadow map discards degrees of freedom that become nonadmissible (or noncoherent) once the
system crosses a horizon barrier.

\begin{definition}[Horizon barrier (shadow-level identification)]
A horizon barrier is an admissibility barrier $\mathcal B_{\mathrm{hor}}\subset\mathcal X$ such that
crossing it corresponds, in the shadow $\mathcal Y$, to formation of a trapped region and
subsequent evaporation in which interior degrees of freedom no longer remain in the admissible
coherent sector as seen by external observables.
\end{definition}

\begin{proposition}[Information loss as projected noninvertibility]
Under the horizon barrier identification, the nonexistence of a global 4D recovery map for the
pre-collapse state from late-time Hawking radiation is an instance of Theorem~\ref{thm:barrier-noninvertibility}.
\end{proposition}

\subsection{Mapping to Page curve and islands}

The island formula program restores an effective Page curve by modifying how entanglement
entropy of radiation is computed, effectively including additional regions (islands) in the
entanglement wedge. This can be interpreted as enlarging or re-encoding the shadow algebra of
observables.

\begin{remark}[Islands as partial inverse / re-encoding]
In the present framework, islands correspond to regime-dependent partial right-inverses on a
\emph{restricted} algebra or coarse-grained sector of $\mathcal Y$, not a global right-inverse of
$T_t$ on all of $\mathcal Y$. That is, islands can restore invertibility for selected observables
within a universality regime, while Theorem~\ref{thm:barrier-noninvertibility} forbids global
invertibility once barrier-crossing fiber identifications occur.
\end{remark}

\begin{remark}[Mainstream patches predicted by the bridge]
Theorem~\ref{thm:barrier-noninvertibility} predicts that any attempt to restore full 4D
unitarity/invertibility must either:
\begin{enumerate}[label=(\roman*)]
\item enlarge the observable algebra/state space (nonlocal encoding), or
\item introduce explicit final-state projections, or
\item rely on regime-dependent effective inversions (islands) that cannot be global.
\end{enumerate}
This aligns structurally with the emergence of islands/replica-wormhole methods and with
final-state projection proposals.
\end{remark}

% ==========================================================
\section{Confrontation with the Page Curve and Island Formula}
% ==========================================================

In this section we confront the present framework directly with the modern ``island''
resolution of the black hole information problem. Our goal is not to re-derive the island
formula or replica trick, but to clarify—at the level of operator structure and
invertibility—what such constructions can and cannot accomplish in light of
Theorem~\ref{thm:barrier-noninvertibility}.

\subsection{What the island formula actually restores}

The island formula modifies the computation of von Neumann entropy for radiation
subsystems by enlarging the region whose entanglement wedge is included in the
entropy calculation. Schematically, one replaces
\[
S(\rho_{\mathrm{rad}}) \;\longrightarrow\;
\min_{\mathcal I}
\left\{
\frac{\mathrm{Area}(\partial \mathcal I)}{4G_N}
+
S_{\mathrm{bulk}}(\rho_{\mathrm{rad}\cup \mathcal I})
\right\}.
\]

Operationally, this corresponds to enlarging the effective algebra of observables used
to compute entropy, so that certain interior degrees of freedom are treated as if they
were encoded in the radiation sector.

From the perspective of the present framework, this amounts to the following:

\begin{itemize}
\item The observable shadow space $\mathcal Y$ is \emph{effectively enlarged or
re-encoded} for a restricted class of observables (those entering the entropy
functional).
\item The projection $P:\mathcal X\to\mathcal Y$ is \emph{not inverted globally}, but
is partially inverted or redefined on a chosen subalgebra.
\end{itemize}

Thus the island prescription restores \emph{effective invertibility} for specific
coarse-grained diagnostics (entanglement entropy), but does not provide a global
right-inverse for the shadow evolution map $T_t=P\circ\Phi_t$.

\subsection{Islands as partial inverses on restricted algebras}

Let $\mathcal A(\mathcal Y)$ denote the full algebra of 4D observables, and let
$\mathcal A_{\mathrm{island}}\subset \mathcal A(\mathcal Y)$ denote the restricted
subalgebra for which the island prescription applies.

In the present language, the island construction defines a map
\[
S_t^{(\mathcal A_{\mathrm{island}})} :
\mathcal A_{\mathrm{island}} \;\to\; \mathcal X
\]
such that
\[
T_t \circ S_t^{(\mathcal A_{\mathrm{island}})}
=
\Id
\quad \text{on } \mathcal A_{\mathrm{island}}.
\]

Crucially, no such map exists on the full algebra $\mathcal A(\mathcal Y)$ once
barrier-crossing noninjectivity has occurred.

\begin{remark}
This distinction is often blurred in informal discussions of ``information recovery.''
Theorem~\ref{thm:barrier-noninvertibility} forbids a global inverse on $\mathcal Y$, but
does not forbid partial inverses on restricted algebras or coarse-grained observables.
\end{remark}

\subsection{Why islands do not restore global unitarity}

Theorem~\ref{thm:barrier-noninvertibility} implies that once the shadow map $T_t$ identifies
distinct upstairs configurations across an admissibility barrier, no measurable global
right-inverse can exist.

Island constructions evade this only by:
\begin{enumerate}[label=(\roman*)]
\item restricting attention to specific observables (entropy functionals), and
\item allowing state-dependent or regime-dependent re-encodings of degrees of freedom.
\end{enumerate}

Neither step restores global invertibility of the shadow dynamics. Instead, islands
implement a \emph{contextual partial reconstruction}—exactly analogous, in structure, to
context-dependent recovery of pre-measurement information in quantum measurement theory.

\subsection{Predictions of the MTT shadow-bridge}

The present framework makes sharp structural predictions that go beyond the island
literature:

\begin{enumerate}[label=(P\arabic*)]
\item \textbf{No protocol can reconstruct full pre-collapse data.} Any procedure claiming
full information recovery must either enlarge the observable algebra beyond $\mathcal Y$
or introduce nonlocal degrees of freedom not present in the original shadow theory.
\item \textbf{Page-curve unitarity is effective, not fundamental.} The recovery of a Page
curve reflects a reorganization of basin measures and partial re-encodings, not restoration
of global invertibility.
\item \textbf{Breakdown occurs at admissibility barriers, not entropy thresholds.} The
fundamental transition is controlled by gap closure / loss of projector regularity, not by
the entanglement entropy reaching a maximum.
\end{enumerate}

These predictions distinguish the MTT shadow-bridge from interpretations that treat the
island formula as evidence for exact 4D unitarity.

\subsection{Relation to final-state projection and nonlocal encodings}

Other proposed resolutions—final-state projection at the singularity, ER=EPR-style
nonlocal encoding, or state-dependent interior operators—can be classified in the same
way. Each introduces, implicitly or explicitly, a modification of the observable map
$P$ or an enlargement of $\mathcal Y$ that allows partial inversion on selected sectors.

The shadow-bridge framework predicts that all such constructions are \emph{patches} for the
same underlying fact: projection across an admissibility barrier is noninvertible, and any
restoration of information must therefore be either partial, contextual, or nonlocal.

\subsection{Summary}

The island formula does not contradict the present framework. Rather, it fits naturally
as a regime-dependent partial inversion of the shadow map on restricted observables. What
it cannot do—by Theorem~\ref{thm:barrier-noninvertibility}—is restore global invertibility
of the 4D shadow dynamics once admissibility-barrier crossing has occurred.

In this sense, black hole information loss and quantum measurement collapse share not only
a common origin, but also a common pattern of attempted resolution: effective recovery on
restricted algebras, without global reconstruction.
% ==========================================================
\section{Born Weights and Hawking Weights from a Single Basin-Measure Functional}
% ==========================================================

We now show that the probability weights appearing in quantum measurement (Born rule)
and the thermal weights appearing in Hawking radiation arise from the same mathematical
object in the MTT framework: a basin-measure functional on the full configuration space,
evaluated across different admissibility barriers.

\subsection{Basin measures in the full configuration space}

Let $(\mathcal X,\mu)$ be the MTT configuration space equipped with its natural invariant
measure $\mu$ (Liouville-type, or stationary measure induced by the deterministic flow
$\Phi_t$). Let $\mathcal B\subset\mathcal X$ be an admissibility barrier, and suppose that
crossing $\mathcal B$ induces a decomposition of the admissible set into disjoint basins
\[
\mathcal A \;\longrightarrow\; \bigsqcup_{i\in I} \mathcal A_i,
\]
where each $\mathcal A_i$ flows under $\Phi_t$ to a distinct coherent attractor or sector.

\begin{definition}[Basin-measure functional]
Define the basin-measure functional
\[
W_i := \frac{\mu(\mathcal A_i)}{\sum_{j\in I} \mu(\mathcal A_j)}.
\]
\end{definition}

By construction, $\{W_i\}$ defines a probability distribution over admissible outcomes
after barrier crossing.

\subsection{Measurement barrier and Born weights}

In the measurement setting, the admissibility barrier $\mathcal B_{\mathrm{meas}}$
corresponds to a change of context or coupling to an apparatus, producing a set of outcome
basins $\{\mathcal A_i\}$ associated with stable pointer states.

\begin{proposition}[Born weights as basin measures (structural identification)]

Under the measurement admissibility barrier $\mathcal B_{\mathrm{meas}}$, the Born weights
$p_i = \|\Pi_i \psi\|^2$ coincide with the basin-measure functional $W_i$ induced by $\mu$,
up to normalization.
This identification relies on the existence and invariance of the upstairs measure $\mu$
and on the correspondence between $\mu$ and squared amplitudes established in earlier
MTT $\to$ QM derivations; it is not re-derived here.

\end{proposition}

\begin{remark}
This identification does not require postulating the Born rule. It follows from the
invariance of $\mu$ under $\Phi_t$ and the fact that projection-induced basin capture
is the only source of stochasticity in the shadow dynamics.
\end{remark}

\subsection{Horizon barrier and Hawking weights}

In the black hole setting, the admissibility barrier $\mathcal B_{\mathrm{hor}}$
corresponds to horizon formation and evaporation, which renders interior degrees of
freedom nonadmissible in the coherent shadow. The admissible exterior sector decomposes
into basins labeled by asymptotic radiation modes.

\begin{proposition}[Hawking weights as basin measures under semiclassical identification]

Under the horizon admissibility barrier $\mathcal B_{\mathrm{hor}}$, the thermal Hawking
weights
\[
p_\omega \propto e^{-\beta \omega}
\]
arise as the basin-measure functional $W_\omega$ induced by $\mu$, evaluated on radiation
basins labeled by asymptotic frequency $\omega$.
Here $\beta$ is identified with the inverse Hawking temperature in the standard
semiclassical regime; no derivation of $\beta$ or of horizon microphysics is attempted
in this work.

\end{proposition}

\begin{remark}
The appearance of a thermal spectrum reflects the fact that the invariant measure $\mu$,
when restricted to the exterior-admissible basins, is weighted by the same exponential
action factors that control stability and admissibility in the coherent sector.
\end{remark}

\subsection{Unified theorem}

We can now state the unifying result.

\begin{theorem}[Unifying structural principle]
\label{thm:born-hawking}
Let $\mathcal B$ be an admissibility barrier in an MTT system, inducing basin decomposition
$\{\mathcal A_i\}$ of the admissible set. Then the probability weights observed in the 4D
shadow dynamics after crossing $\mathcal B$ are given universally by the basin-measure
functional $W_i$.

In particular:
\begin{enumerate}[label=(\roman*)]
\item for $\mathcal B=\mathcal B_{\mathrm{meas}}$, $W_i$ yields Born probabilities;
\item for $\mathcal B=\mathcal B_{\mathrm{hor}}$, $W_i$ yields Hawking thermal weights.
\end{enumerate}
\end{theorem}

\begin{corollary}[No fundamental distinction between quantum and thermal probabilities]
In the MTT framework, quantum measurement probabilities and black hole thermal radiation
probabilities are not fundamentally different objects. Both are shadows of the same
upstairs basin-measure functional evaluated across different admissibility barriers.
\end{corollary}

\begin{remark}[Why this is not ``everything is thermal'']
The distinction between ``quantum'' and ``thermal'' statistics lies not in the measure
functional itself, but in the geometry of the barrier and the labeling of the basins.
Measurement barriers produce discrete outcome basins; horizon barriers produce
continuum-labeled radiation basins.
\end{remark}
% ==========================================================
\section{Quantum Measurement Collapse as Partial Inversion on Restricted Algebras}
% ==========================================================

We now complete the shadow-bridge by showing that quantum measurement collapse is
structurally identical to the island construction discussed in the black hole setting.
In both cases, apparent information loss is resolved by a partial inversion of the
shadow map on a restricted algebra of observables.

\subsection{Measurement update as a noninvertible shadow map}

Let $\mathcal A(\mathcal Y)$ denote the algebra of 4D observables. A measurement context
$C$ (specified by an instrument or POVM $\{E_i\}$) induces a restriction of admissibility
and hence a change of the effective observable map $P$.

After a measurement, the shadow dynamics is described by
\[
\rho \;\longrightarrow\; \rho_i
= \frac{\sqrt{E_i}\,\rho\,\sqrt{E_i}}{\Tr(E_i\rho)},
\]
which is manifestly noninvertible on $\mathcal A(\mathcal Y)$.

\subsection{Restricted algebra and partial inverse}

Define the post-measurement algebra
\[
\mathcal A_i := \{ A \in \mathcal A(\mathcal Y) \mid [A,E_i]=0 \},
\]
the algebra of observables compatible with outcome $i$.

\begin{proposition}[Partial inversion after measurement]
For each outcome $i$, there exists a map
\[
S_i : \mathcal A_i \;\to\; \mathcal X
\]
such that
\[
P \circ \Phi_t \circ S_i = \Id
\quad \text{on } \mathcal A_i.
\]
No such map exists on the full algebra $\mathcal A(\mathcal Y)$.
\end{proposition}

\begin{remark}
This is the measurement analogue of the island construction: invertibility is restored
only on a restricted algebra selected by the context and outcome.
\end{remark}

\subsection{Collapse vs.\ Everett vs.\ islands}

From this perspective:
\begin{itemize}
\item ``Collapse'' is the statement that the shadow map is noninvertible on the full algebra.
\item ``Many worlds'' corresponds to keeping the full upstairs state in $\mathcal X$ and
refusing to project.
\item ``Decoherence'' corresponds to suppressing off-diagonal sectors without selecting
a basin.
\item ``POVM update'' is a partial inversion on $\mathcal A_i$, analogous to an island.
\end{itemize}

\subsection{Unified picture}

\begin{theorem}[Measurement--island equivalence]
\label{thm:measurement-island}
Quantum measurement collapse and black hole island constructions are instances of the
same structural mechanism: partial inversion of a noninvertible shadow map on a restricted
observable algebra, following admissibility-barrier crossing.
\end{theorem}

\begin{remark}
The difference between measurement and black holes lies in the physical realization of
the barrier and in the algebra selected for partial inversion, not in the underlying
mechanism.
\end{remark}
% ==========================================================
\section{Direct Comparison with Everettian Interpretations}
% ==========================================================

Everettian (``many-worlds'') interpretations propose that the appearance of collapse is an
illusion arising from unitary evolution and branching of the universal wavefunction, with
no fundamental projection or loss of information. In this section we confront that claim
directly within the shadow-bridge framework. This distinction should be understood as a theory-choice criterion about what counts
as an autonomous effective description of 4D physics, not as an empirical refutation
of Everettian interpretations.


\subsection{The Everettian move}

The Everettian position may be summarized as follows:
\begin{enumerate}[label=(E\arabic*)]
\item The fundamental state (wavefunction) evolves unitarily at all times.
\item Measurement corresponds to entanglement and branching, not projection.
\item Apparent collapse reflects an observer's restriction to a branch, not a physical
process.
\end{enumerate}

From the present perspective, this corresponds to \emph{refusing to apply the projection}
$P=I\circ\Pi$ and insisting that the physically relevant description remains the upstairs
state in $\mathcal X$ (or its Hilbert-space analogue) at all times.

\subsection{What Everett must deny}

Theorem~\ref{thm:barrier-noninvertibility} isolates the precise point of disagreement.
The Everettian position requires that:
\begin{itemize}
\item the full upstairs state remains physically meaningful at the 4D level, and
\item no admissibility barrier produces noninvertibility of the observable map.
\end{itemize}

In other words, Everett denies that $P$ is a physically enforced map, treating it instead
as an optional epistemic restriction.

\subsection{Structural consequences of refusing projection}

If one refuses the projection $P$ as physically operative, two consequences follow:

\begin{enumerate}[label=(C\arabic*)]
\item The effective 4D theory is no longer complete: observers must be described as parts
of the upstairs configuration space $\mathcal X$, with no autonomous 4D description.
\item All apparent irreversibility (collapse, thermalization, horizon loss) must be
reconstructed from branch-relative bookkeeping rather than from dynamical selection.
\end{enumerate}

This position is internally consistent, but it abandons the goal of deriving a closed,
autonomous 4D effective theory.

\subsection{Everett and black holes}

The same structural choice appears in the black hole context. An Everettian stance would
treat the full interior--exterior entangled state as physically real and deny any loss of
information, regardless of the inability of 4D observers to reconstruct it.

In the shadow-bridge framework, this corresponds to refusing the projection across the
horizon barrier and insisting that interior degrees of freedom remain part of the physical
state, even when they are nonadmissible in the coherent sector.

\subsection{Why Everett reintroduces shadows implicitly}

In practice, Everettian treatments inevitably reintroduce shadow-like structures:
\begin{itemize}
\item Branch weights obey the Born rule.
\item Observers experience definite outcomes.
\item Decoherence defines effective pointer sectors.
\end{itemize}

These are precisely the signatures of basin selection and partial inversion on restricted
algebras described earlier. From the shadow-bridge perspective, Everett retains the
upstairs state but relies on emergent shadow structure to recover phenomenology.

\begin{remark}[Everett as a refusal, not a refutation]
The present framework does not refute Everettian interpretations. Rather, it clarifies
their commitment: Everett denies that projection-induced noninvertibility is physically
operative. In exchange, it gives up a closed 4D effective description and must treat all
observed definiteness as branch-relative.
\end{remark}

\subsection{Shadow-bridge position}

The shadow-bridge framework makes the opposite choice:
\begin{itemize}
\item The upstairs evolution is deterministic and invertible.
\item The projection $P$ is physically enforced by admissibility and stability.
\item Noninvertibility of the 4D shadow is real and generative, not illusory.
\end{itemize}

From this standpoint, collapse and black hole information loss are not interpretational
artifacts but structural features of effective physics under projection.

\begin{theorem}[Everett--shadow dichotomy]
\label{thm:everett-shadow}
Everettian interpretations and the shadow-bridge framework are distinguished by a single
structural choice: whether the projection $P$ is treated as physically operative. If $P$
is operative, barrier-crossing noninvertibility is unavoidable and collapse/information
loss are real in the 4D shadow. If $P$ is refused, 4D physics is not closed and must be
understood as branch-relative.
\end{theorem}

\begin{remark}
This dichotomy is not empirical at the level of the upstairs theory, but it is decisive at
the level of what counts as an autonomous effective description of 4D physics.
\end{remark}

\section{Outlook}

This paper isolates the common mathematical core of two crises: black hole information loss and
quantum measurement collapse. The next step is to instantiate the barrier structure in concrete
MTT models: identify the spectral-gap/contractivity conditions corresponding to measurement
context changes and to horizon formation, and compute the induced shadow entropy diagnostics
in each case. A direct comparison with Page-curve/island formula computations would then
locate precisely which observables admit partial inversion and where the coherent regime must
fail.
\begin{thebibliography}{99}

\bibitem{MTT-Admissibility}
P.~Nero,
\emph{Modal Triplet Theory: Admissibility, Encodings, and the Structure of Physical Description},
Zenodo preprint, January 2026.
\url{https://doi.org/10.5281/zenodo.18255621}

\bibitem{MTT-Foundation}
P.~Nero,
\emph{Modal Triplet Theory: Foundation},
Zenodo preprint, September 2025.
\url{https://doi.org/10.5281/zenodo.16949762}

\bibitem{MTT-FixedPoints}
P.~Nero,
\emph{Fixed Points I--VI: Complete Coherence Spine},
Zenodo preprints, August 2025.
\url{https://doi.org/10.5281/zenodo.16948748}

\bibitem{MTT-ProjectionAdmissibility}
P.~Nero,
\emph{The Projection--Admissibility Principle: Structural Constraints on Effective Physical Description},
Zenodo preprint, January 2026.
\url{https://doi.org/10.5281/zenodo.18255838}

\bibitem{MTT-Closure}
P.~Nero,
\emph{Closure and Inevitability in Modal Triplet Theory},
Zenodo preprint, January 2026.
\url{https://doi.org/10.5281/zenodo.18255510}

\bibitem{MTT-CoherenceCapacity}
P.~Nero,
\emph{Coherence Capacity as the Fundamental Resource of Effective Physics},
Zenodo preprint, January 2026.
\url{https://doi.org/10.5281/zenodo.18255905}

\bibitem{MTT-CoherenceDynamics}
P.~Nero,
\emph{Dynamics of Coherence Capacity: Transport, Concentration, and Exhaustion},
Zenodo preprint, January 2026.
\url{https://doi.org/10.5281/zenodo.18256048}

\bibitem{MTT-QM}
P.~Nero,
\emph{Modal Triplet Theory: From MTT to Quantum Mechanics},
Zenodo preprint, September 2025.
\url{https://doi.org/10.5281/zenodo.17074246}

\bibitem{MTT-QFT}
P.~Nero,
\emph{From MTT to Quantum Field Theory},
Zenodo preprint, 2025.
\url{https://doi.org/10.5281/zenodo.17068816}

\bibitem{MTT-GR}
P.~Nero,
\emph{Modal Triplet Theory: From MTT to General Relativity},
Zenodo preprint, October 2025.
\url{https://doi.org/10.5281/zenodo.16950597}

\bibitem{MTT-UVQG}
P.~Nero,
\emph{Modal Triplet Theory: From MTT to a UV-Finite, Unitary Quantum Gravity},
Zenodo preprint, 2025.
\url{https://doi.org/10.5281/zenodo.17077671}

\bibitem{MTT-Measurement}
P.~Nero,
\emph{Measurement as Disturbance and Stabilization in Modal Triplet Theory},
Zenodo preprint, 2025.
\url{https://doi.org/10.5281/zenodo.17177404}

\bibitem{MTT-Irreversibility}
P.~Nero,
\emph{Projection, Probability, and Irreversibility: Shadow Bridges Between Measurement, Black Holes, and Cosmology in Modal Triplet Theory},
Zenodo preprint, January 2026.
\url{https://doi.org/10.5281/zenodo.18256408}

\bibitem{MTT-Bell-Beables}
P.~Nero,
\emph{Modal Fixed Points, Bell’s Beables, and the Limits of Factorization},
Zenodo preprint, 2025.
\url{https://doi.org/10.5281/zenodo.17076300}

\bibitem{MTT-Bell-Temporal}
P.~Nero,
\emph{Temporal Bell Inequalities and Global Consistency in Modal Triplet Theory},
Zenodo preprint, August 2025.
\url{https://doi.org/10.5281/zenodo.18208884}

\bibitem{MTT-Stochastic}
P.~Nero,
\emph{From Modal Triplet Theory to Indivisible Stochastic Processes: A First-Principles, Fully Rigorous Derivation},
Zenodo preprint, January 2026.
\url{https://doi.org/10.5281/zenodo.18254862}

\end{thebibliography}
\end{document}
